\documentclass[10pt,a4paper]{amsart}
\usepackage[margin=19mm]{geometry}
\usepackage{amsmath,amssymb,mathtools}
\usepackage[scr=rsfs,cal=euler]{mathalfa}
\usepackage{xfrac}
\usepackage[unicode,hidelinks,bookmarks=false]{hyperref}
\newcommand{\ie}{i.e.\ }
\newcommand{\cf}{cf.\ }
\newcommand{\resp}{respectively\ }
\newcommand{\Subsection}{\subsection}
\providecommand{\nobreakdash}{\nobreak\hbox{-}}
\setlength{\emergencystretch}{2em}
\multlinegap0pt
\usepackage{mathrsfs}
\usepackage{amstext}
\numberwithin{equation}{section}
\numberwithin{figure}{section}
\newtheorem{thm}{\protect\theoremname}
\newtheorem{conjecture}[thm]{\protect\conjecturename}
\newtheorem{cor}[thm]{\protect\corollaryname}
\newtheorem{defn}[thm]{\protect\definitionname}
\newtheorem{prop}[thm]{\protect\propositionname}
\providecommand{\conjecturename}{Conjecture}
\providecommand{\corollaryname}{Corollary}
\providecommand{\definitionname}{Definition}
\providecommand{\propositionname}{Proposition}
\providecommand{\theoremname}{Theorem}
\begin{document}
\title[Averages of diagonal Elliott-Halberstam problem twisted by M]{Averages of diagonal Elliott-Halberstam problem twisted by M\"obius function with Sobolev and H\"older-Zygmund weights}
\author{Marco Cantarini}
\date{}
\maketitle
\noindent\emph{Source and licence.} Averages of diagonal Elliott-Halberstam problem twisted by M\"obius function with Sobolev and H\"older-Zygmund weights. Version arXiv:2607.09110v1 (CC BY 4.0). This material is distributed under the Creative Commons Attribution 4.0 International licence, http://creativecommons.org/licenses/by/4.0/, as recorded in the arXiv OAI licence metadata for this exact version and shown on the abstract page https://arxiv.org/abs/2607.09110v1. Full rights notice: licenses/arxiv-2607.09110-CC-BY-4.0.txt.\par\smallskip
\noindent\textit{Editorial scope.} Counted unit: the original \texttt{thm} environment \texttt{thm:dobule series =00005Cchi} at source lines 903--937 together with its complete proof at lines 939--1035; the delivered document keeps source lines 587--1044 so that every referenced label and every citation resolves inside it. Native editable source is the paper's own concatenated TeX; the AMS wrapper is editorial formatting only. No publisher class, upstream script or unrelated artwork is redistributed.
\section{preliminary results, definitions and main tools}

In this section we have to recall several facts and prove some preliminary
results. For this reason, we divide this section further subsections.

\subsection{The main tools}

One of the main tools used to show our theorems is an identity proved
in \cite{CGZ2025}. It can be viewed as a two-dimensional Abel summation
formula. Versions of two-dimensional Abel summation do appear in the
literature (see, e.g, \cite{KT2006} or \cite{BP2018}), but the presented
formulation makes explicit the link between a weighted average involving
two arithmetic functions and the Laplace convolution of the corresponding
unweighted averages taken separately. This point of view is particularly
useful, since it allows us to combine explicit formulae for the individual
averages, whenever available, and thereby derive an explicit formula
for the original weighted average. To the best of our knowledge, Abel
summation has not previously been formulated in exactly this way in
the literature.
\begin{thm}
\label{thm:MAINCGZ}Let $g_{1},g_{2}$ be arithmetical functions,
$\eta\in\mathbb{R}^{+}$, $f:\mathbb{R}\rightarrow\mathbb{C}$, and
assume that:

1) $f$ has its support in $\left[\alpha,\beta\right),\,0\leq\alpha<\beta$,
and $\beta\in\mathbb{R}$

2) $f\in C^{1}\left(\alpha,\beta\right)$.

3) $f^{\prime}\in AC\left(\alpha,\beta\right)$.

Then, if 
\[
G_{j}\left(x\right):=\begin{cases}
\sum_{n\leq x}g_{j}\left(n\right), & x>0\\
0, & \text{otherwise}
\end{cases}
\]
for $j=1,2$, we get
\[
\sum_{\eta\alpha<n\leq\eta\beta}\sum_{m\leq\eta\beta-n}g_{2}\left(m\right)g_{1}\left(n\right)f\left(\frac{m+n}{\eta}\right)=G_{2}\left(\eta\alpha\right)\int_{\alpha}^{\beta}G_{1}\left(\eta v-\eta\alpha\right)f^{\prime}\left(v\right)dv
\]
\[
+\frac{1}{\eta}\int_{\alpha}^{\beta}f^{\prime\prime}\left(w\right)\int_{\eta\alpha}^{\eta w}G_{2}\left(s\right)G_{1}\left(\eta w-s\right)dsdw.
\]
\end{thm}

Observe that, essentially, we require that $f_{\vert\left(\alpha,\beta\right)}$
belongs to the Sobolev spaces $W^{2,1}\left(\alpha,\beta\right)$.
We recall the definition, just for completeness: if $I$ is an open
interval, then
\[
W^{1,p}\left(I\right):=\left\{ f\in L^{p}\left(I\right):\,\exists g\in L^{p}\left(I\right):\int_{I}fh^{\prime}=-\int_{I}gh,\,\forall h\in C_{c}^{1}\left(I\right)\right\} 
\]
where $C_{c}^{1}\left(I\right)$ are the function $C^{1}\left(I\right)$
with compact support in $I$, and
\[
W^{m,p}\left(I\right):=\left\{ f\in W^{m-1,p}\left(I\right):f^{\prime}\in W^{m-1,p}\left(I\right)\right\} ,\,m\in\mathbb{N},\,m\ge2
\]
(see, e.g., \cite{B2011}, chapt. $8)$. We recall that is possible
to extend the results of the previous theorem also in the case $\beta=+\infty$
(see \cite{CGZ2025}). Note that the hypotheses of the Theorem imply
that $f\left(\alpha^{+}\right),f^{\prime}\left(\alpha^{+}\right)$
exists and are finite (but not necessary equal to $0$) and $f\left(\beta^{-}\right)=f^{\prime}\left(\beta^{-}\right)=0$. 

In order to work with more general weights, we propose now a discrete
version of the previous theorem. Firstly, we recall the definition
of $r$-th order forward difference, with $r\in\mathbb{N}^{+}$, as
\[
\Delta_{h}^{r}\left(f,x\right):=\sum_{k=0}^{r}\binom{r}{k}\left(-1\right)^{r-k}f\left(x+hk\right),\,x,h\in\mathbb{R},
\]
\[
\Delta_{h}\left(f,x\right):=\Delta_{h}^{1}\left(f,x\right)
\]
and the property
\begin{equation}
\Delta_{h}^{r}=\Delta_{h}\left[\Delta_{h}^{r-1}\right]\label{eq:prod fow diff}
\end{equation}
(see. e.g., \cite{DVL1993}, chapt. $7$).
\begin{thm}
\label{thm:MainDiscrete}Let $g_{1},g_{2}$ be arithmetical functions,
let $f:\mathbb{R}\rightarrow\mathbb{C}$ and assume that:

1) f has its support in $[\alpha,\beta)$, $0\leq\alpha<\beta$$.$ 

2) $G_{1}(1)=G_{1}(0)=G_{2}(0)=0.$

Then, taking $\eta\in\mathbb{R}^{+}$ such that $\eta\beta\in\mathbb{N}^{+},\,\eta\alpha\in\left\{ 0,1\right\} $
we have
\[
\sum_{\eta\alpha<n\leq\eta\beta}\sum_{m\leq\eta\beta-n}g_{2}\left(m\right)g_{1}\left(n\right)f\left(\frac{m+n}{\eta}\right)=\sum_{k=2}^{\eta\beta-1}\left(\sum_{u=0}^{k}G_{1}\left(u\right)G_{2}\left(k-u\right)\right)\Delta_{1/\eta}^{2}\left(f,\frac{k}{\eta}\right).
\]
\end{thm}

\begin{proof}
By classical summation by parts and since $f(\beta)=f\left(\beta+\frac{1}{\eta}\right)=0$,
because $f$ has its support in $[\alpha,\beta)$, we have
\[
\sum_{m\leq\eta\beta-n}g_{2}\left(m\right)f\left(\frac{m+n}{\eta}\right)=-\sum_{m\leq\eta\beta-n-1}G_{2}(m)\left[f\left(\frac{m+n+1}{\eta}\right)-f\left(\frac{m+n}{\eta}\right)\right]
\]
and so defining 
\[
C(n):=-\sum_{m\leq\eta\beta-n-1}G_{2}(m)\left[f\left(\frac{m+n+1}{\eta}\right)-f\left(\frac{m+n}{\eta}\right)\right]
\]
we have, again by partial summation
\[
\sum_{\eta\alpha<n\leq\eta\beta}g_{1}\left(n\right)\sum_{m\leq\eta\beta-n}g_{2}\left(m\right)f\left(\frac{m+n}{\eta}\right)=G_{1}\left(\eta\beta\right)C\left(\eta\beta\right)-G_{1}\left(\eta\alpha\right)C\left(\eta\alpha+1\right)
\]
\[
-\sum_{\eta\alpha<n\leq\eta\beta-1}G_{1}\left(n\right)\left(C\left(n+1\right)-C\left(n\right)\right).
\]
Now clearly
\[
C\left(\eta\beta\right)=G_{1}\left(\eta\alpha\right)=0
\]
hence
\[
\sum_{\eta\alpha<n\leq\eta\beta}g_{1}\left(n\right)\sum_{m\leq\eta\beta-n}g_{2}\left(m\right)f\left(\frac{m+n}{\eta}\right)=-\sum_{\eta\alpha<n\leq\eta\beta-1}G_{1}\left(n\right)\left(C\left(n+1\right)-C\left(n\right)\right)
\]
\[
=\sum_{\eta\alpha<n\leq\eta\beta-1}G_{1}\left(n\right)\sum_{m\leq\eta\beta-n-2}G_{2}(m)\left[f\left(\frac{m+n+2}{\eta}\right)-f\left(\frac{m+n+1}{\eta}\right)\right]
\]
\[
-\sum_{\eta\alpha<n\leq\eta\beta-1}G_{1}\left(n\right)\sum_{m\leq\eta\beta-n-1}G_{2}(m)\left[f\left(\frac{m+n+1}{\eta}\right)-f\left(\frac{m+n}{\eta}\right)\right]
\]
\[
=\sum_{\eta\alpha<n\leq\eta\beta-1}G_{1}\left(n\right)\sum_{m\leq\eta\beta-n-1}G_{2}(m)\Delta_{1/\eta}^{2}\left(f,\frac{m+n}{\eta}\right)
\]
then, taking $k=m+n$ and using condition $2)$, we obtain
\[
\sum_{\eta\alpha<n\leq\eta\beta}g_{1}\left(n\right)\sum_{m\leq\eta\beta-n}g_{2}\left(m\right)f\left(\frac{m+n}{\eta}\right)=\sum_{\eta\alpha<k\leq\eta\beta-1}\Delta_{1/\eta}^{2}\left(f,\frac{k}{\eta}\right)\sum_{u=0}^{k}\left(G_{1}(u)G_{2}(k-u)\right)
\]
as wanted.
\end{proof}
%
The previous result can be proved in a more general settings, but
paying the price of having an identity with a more complicated form.
\begin{cor}
\label{cor:discrete}Under the previous hypotheses, we have
\[
\sum_{\eta\alpha<n\leq\eta\beta}g_{1}\left(n\right)\sum_{m\leq\eta\beta-n}g_{2}\left(m\right)f\left(\frac{m+n}{\eta}\right)=\sum_{\eta\alpha<k\leq\eta\beta-1}\Delta_{1/\eta}^{2}\left(f,\frac{k}{\eta}\right)\int_{0}^{k}G_{1}(s)G_{2}(k-s)ds
\]
\[
+\sum_{\eta\alpha<k\leq\eta\beta-1}\Delta_{1/\eta}^{2}\left(f,\frac{k}{\eta}\right)\sum_{u=0}^{k}G_{1}(u)g_{2}(k-u).
\]
\end{cor}

\begin{proof}
It is enough to observe that
\[
\int_{0}^{k}G_{1}(s)G_{2}(k-s)ds=\sum_{u=0}^{k-1}\int_{u}^{u+1}G_{1}(s)G_{2}(k-s)ds=\sum_{u=0}^{k-1}G_{1}(u)G_{2}(k-u-1)
\]
hence
\[
\sum_{u=0}^{k}G_{1}(u)G_{2}(k-u)-\sum_{u=0}^{k-1}G_{1}(u)G_{2}(k-u-1)
\]
\[
=\sum_{u=0}^{k-1}G_{1}(u)\left[G_{2}(k-u)-G_{2}(k-u-1)\right]
\]
\[
=\sum_{u=0}^{k-1}G_{1}(u)g_{2}\left(k-u\right).
\]
\end{proof}

\subsection{H\"older-Zygmund spaces}

As mentioned in the introduction, we are interested in weights defined
by functions belonging to suitable function spaces. In the previous
subsection we considered the case in which the restriction $f_{\vert\left(\alpha,\beta\right)}\in W^{2,1}\left(\alpha,\beta\right)$.
In this subsection, we turn to the second setting under consideration,
namely H\"older-Zygmund spaces. We take this class because the previous
theorems show that we need an effective way to control finite differences
of the weight. We therefore recall the definition of these spaces
and some of their basic properties. Everything will be formulated
on $\mathbb{R},$ since our functions are defined on $\mathbb{R}$,
even when their support is contained in an interval; this is important
because the behavior near the boundary is different.

\begin{defn}
(H\"older-Zygmund space) Let $\delta>0$. We define the H\"older-Zygmund
space as
\[
\mathcal{C}^{\delta}\left(\mathbb{R}\right):=\left\{ f\in L^{\infty}\left(\mathbb{R}\right)\cap C^{0}\left(\mathbb{R}\right):\left\Vert f\right\Vert _{\mathcal{C}^{\delta}\left(\mathbb{R}\right)}<+\infty\right\} 
\]
where
\[
\left\Vert f\right\Vert _{\mathcal{C}^{\delta}\left(\mathbb{R}\right)}:=\left\Vert f\right\Vert _{\infty}+\sup_{x,h\in\mathbb{R},\,h\neq0}\frac{\left|\Delta_{h}^{1}\left(f,x\right)\right|}{h^{\delta}}
\]
if $0<\delta<1$,
\[
\left\Vert f\right\Vert _{\mathcal{C}^{1}\left(\mathbb{R}\right)}:=\left\Vert f\right\Vert _{\infty}+\sup_{x,h\in\mathbb{R},\,h\neq0}\frac{\left|\Delta_{h}^{2}\left(f,x\right)\right|}{h}
\]
if $\delta=1$. If $\delta>1$, say $m<\delta\leq m+1,\,m\in\mathbb{N}$,
then we define the H\"older-Zygmund space as
\[
\mathcal{C}^{\delta}\left(\mathbb{R}\right):=\left\{ f\in L^{\infty}\left(\mathbb{R}\right)\cap C^{m}\left(\mathbb{R}\right):\left\Vert f\right\Vert _{\mathcal{C}^{\delta}\left(\mathbb{R}\right)}<+\infty\right\} 
\]
where the norm is defined recursively as
\[
\left\Vert f\right\Vert _{\mathcal{C}^{\delta}\left(\mathbb{R}\right)}=\left\Vert f\right\Vert _{\mathcal{C}^{\delta-1}\left(\mathbb{R}\right)}+\left\Vert f^{\prime}\right\Vert _{\mathcal{C}^{\delta-1}\left(\mathbb{R}\right)}.
\]
\end{defn}

(see, e.g, \cite{R2022}) Note that $\mathcal{C}^{1}\left(\mathbb{R}\right)$
is also called Zygmund spaces, and if $\delta\notin\mathbb{N}$, $\mathcal{C}^{\delta}\left(\mathbb{R}\right)$
is equivalent to the classical H\"older spaces. Moreover, we have
the following result.
\begin{thm}
\label{thm:finitediff bound}If $f\in L^{\infty}\left(\mathbb{R}\right)\cap C^{0}\left(\mathbb{R}\right)$
and $0<\delta<n,\,n\in\mathbb{N}$, then $f\in\mathcal{C}^{\delta}\left(\mathbb{R}\right)\Leftrightarrow\left|\Delta_{h}^{n}\left(f,x\right)\right|\leq C\left|h\right|^{\delta}$
for all $x,h\in\mathbb{R}$.
\end{thm}

For a reference, see \cite{K1983}, Theorem 6.1. We also recall that
the following useful relations for function in Zygumd and H\"older--Zygmund
class about the definite difference $\Delta_{h}^{1}\left(f,x\right)$;
indeed, we have
\begin{equation}
f\in\mathcal{C}^{1}\left(\mathbb{R}\right)\Rightarrow\left|\Delta_{h}^{1}\left(f,x\right)\right|\leq M\left|h\log\left(\frac{1}{h}\right)\right|,\,M>0\label{eq:LogLip}
\end{equation}
for all $x,h\in\mathbb{R},\,h\neq0$ (see \cite{C1995}, Prop. 2.3.7)
and, if $\delta>1$
\[
f\in\mathcal{C}^{\delta}\left(\mathbb{R}\right)\Rightarrow\left|\Delta_{h}^{1}\left(f,x\right)\right|\leq M\left|h\right|,\,M>0
\]
that is, $f$ verifies the classical Lipschitz condition, for all
$x,h\in\mathbb{R}$, since
\[
\mathcal{C}^{\delta}\left(\mathbb{R}\right)\hookrightarrow\text{Lip}\left(\mathbb{R}\right)
\]
if $\delta>1$, where 
\[
\text{Lip}\left(\mathbb{R}\right):=\left\{ f\in L^{\infty}\left(\mathbb{R}\right):\left\Vert f\right\Vert _{\text{Lip}\left(\mathbb{R}\right)}<+\infty\right\} 
\]
and
\[
\left\Vert f\right\Vert _{\text{Lip}\left(\mathbb{R}\right)}=\left\Vert f\right\Vert _{\infty}+\sup_{x,h\in\mathbb{R},h\neq0}\frac{\left|\Delta_{h}^{1}\left(f,x\right)\right|}{\left|h\right|}
\]
(see again \cite{R2022}). 

\subsection{Estimate for ratio of Gamma function $\Gamma(x)$ and other results}

It is a standard fact that, when one considers a discrete convolution
of functions admitting an explicit formula involving series over zeros,
the corresponding explicit formula for the convolution is related,
in one way or another, to a sum involving a convolution of those zeros.
Very often, this leads to the need to control series involving the
Euler beta function or ratios of gamma functions. This phenomenon
also arises in our work, so it is convenient to record the following
proposition.
\begin{prop}
\label{prop:gamma est}Let $x,y\in\mathbb{R}$. Then
\[
\left|\frac{\Gamma\left(\frac{1}{2}+ix\right)\Gamma\left(\frac{1}{2}+iy\right)}{\Gamma\left(2+i(x+y)\right)}\right|\ll\begin{cases}
\frac{e^{\pi\left(\left|x+y\right|-\left|x\right|-\left|y\right|\right)/2}}{\left|x+y\right|^{3/2}}, & \left|x+y\right|>1\\
e^{-\pi\left(\left|x\right|+\left|y\right|\right)/2}, & \left|x+y\right|\leq1,
\end{cases}
\]
\end{prop}

\begin{proof}
From the well-known relations
\[
\left|\Gamma\left(\frac{1}{2}+ix\right)\right|^{2}=\frac{\pi}{\cosh\left(\pi x\right)},\,\left|\Gamma\left(2+i(x+y)\right)\right|^{2}=\left(1+\left(x+y\right)^{2}\right)\frac{\pi\left|x+y\right|}{\sinh\left(\pi\left|x+y\right|\right)}
\]
(see \cite{O2010}, $5.4.3$ and $5.4.4$), then
\[
\left|\frac{\Gamma\left(\frac{1}{2}+ix\right)\Gamma\left(\frac{1}{2}+iy\right)}{\Gamma\left(2+i(x+y)\right)}\right|^{2}=\pi\frac{\sinh\left(\pi\left|x+y\right|\right)}{\left(1+\left(x+y\right)^{2}\right)\left|x+y\right|\cosh\left(\pi x\right)\cosh\left(\pi y\right)}.
\]
Assume $\left|x+y\right|>1.$ Then, since
\[
\sinh\left(\pi\left|s\right|\right)\leq\frac{e^{\pi\left|s\right|}}{2},\cosh\left(\pi s\right)\geq\frac{e^{\pi\left|s\right|}}{2}
\]
valid for all $s\in\mathbb{R},$ we get
\[
\left|\frac{\Gamma\left(\frac{1}{2}+ix\right)\Gamma\left(\frac{1}{2}+iy\right)}{\Gamma\left(2+i(x+y)\right)}\right|\ll\frac{e^{\pi\left(\left|x+y\right|-\left|x\right|-\left|y\right|\right)/2}}{\left|x+y\right|^{3/2}}.
\]
If $\left|x+y\right|\leq1$, since
\[
\frac{\sinh\left(\pi\left|s\right|\right)}{\pi\left|s\right|}\ll1
\]
then we have
\[
\left|\frac{\Gamma\left(\frac{1}{2}+ix\right)\Gamma\left(\frac{1}{2}+iy\right)}{\Gamma\left(2+i(x+y)\right)}\right|\ll e^{-\pi\left(\left|x\right|+\left|y\right|\right)/2}.
\]
\end{proof}
%
Convergence of double series on the non-trivial zeros is a delicate
matter; for this reason, in \cite{CGZ2026} , we prove a theorem about
the convergence of a double series over the non-trival zeros of the
Riemann zeta function $\zeta(s)$ and a suitable function $f(\rho)$
that verifies some hypotheses. We need a similar results but now for
the non-trivial zeros of $L\left(s,\chi\right)$. We first recall
the following result.
\begin{prop}
Let $T\geq4$ be a real number and $N\left(T\right)$ be the number
of non-trivial zeros $\rho=\sigma+i\gamma$ of $\zeta\left(s\right)$
in the rectangle $0<\sigma<1$, $0<\gamma<T$. Moreover, let $\chi$
a primitive character $\mod q$ with $q\in\mathbb{N}^{+},q>1$ and
let $N\left(T,\chi\right)$ be the number of non-trivial zeros $\rho_{\chi}=\sigma_{\chi}+i\gamma_{\chi}$
in the rectangle $0<\sigma_{\chi}<1,0\leq\gamma_{\chi}\leq T$. Then,
we have
\begin{equation}
N\left(T\right)=\frac{T}{2\pi}\log\left(\frac{T}{2\pi}\right)-\frac{T}{2\pi}+O\left(\log\left(T\right)\right),\label{eq:N(T)}
\end{equation}
\begin{equation}
N\left(T,\chi\right)=\frac{T}{2\pi}\log\left(\frac{qT}{2\pi}\right)-\frac{T}{2\pi}+O\left(\log\left(qT\right)\right).\label{eq:N(T,=00005Cchi)}
\end{equation}
\end{prop}

See, e.g, \cite{DAVEN}, chapt. 15 and 16 or \cite{MV2006}, chapt.
14. Clearly, from the previous formulae we may derive also the cases
$\left|\gamma\right|\leq T,\left|\gamma_{\chi}\right|\leq T$, due
to the symmetry of the zeros of $\zeta\left(s\right)$ and from the
fact that the number of zeros of $L\left(s,\chi\right)$ with $-T\leq\gamma_{\chi}\leq0$
is $N\left(T,\overline{\chi}\right)$.

\begin{thm}
\label{thm:dobule series =00005Cchi}Assume the generalized Riemann
hypothesis, fix $q_{1},q_{2}\in\mathbb{N}^{+}$ and $\chi_{1}\mod q_{1},\,\chi_{2}\mod q_{2}$
two Dirichlet character. Let $\rho_{\chi_{j}}=\frac{1}{2}+i\gamma_{\chi_{j}}$
runs over the non-trivial zeros of the $L$- function $L(s,\chi_{j})$,
$j=1,2$. Let 
\[
Z_{0}^{+}\left(\chi_{j}\right):=\left\{ \rho_{\chi_{j}}:L\left(\rho_{\chi_{j}},\chi_{j}\right)=0,\,\gamma_{\chi_{j}}\geq0\right\} ,
\]
\[
Z^{+}\left(\chi_{j}\right):=\left\{ \rho_{\chi_{j}}:L\left(\rho_{\chi_{j}},\chi_{j}\right)=0,\,\gamma_{\chi_{j}}>0\right\} 
\]
and assume that with the notation
\[
\sum_{\rho_{\chi_{j}}\in Z_{0}^{+}\left(\chi_{j}\right)},\,\sum_{\rho_{\chi_{j}}\in Z^{+}\left(\chi_{j}\right)}
\]
we intend that we are taking $\rho_{\chi_{j}}\in Z_{0}^{+}\left(\chi_{j}\right),\,\rho_{\chi_{j}}\in Z^{+}\left(\chi_{j}\right)$
and summing it considering their multiplicity, for $j=1,2$. Let $f\left(z\right),\,g\left(z\right)$
be complex functions defined on the whole set of the non-trivial zeros
of $L\left(s,\chi_{1}\right)$, $L\left(s,\overline{\chi}_{1}\right)$
and $L\left(s,\chi_{2}\right)$, $L\left(s,\overline{\chi}_{2}\right)$,
respectively. Moreover, assume
\[
\sum_{\underset{{\scriptstyle \gamma_{\chi_{1}}\leq T}}{\rho_{\chi_{1}}\in Z_{0}^{+}\left(\chi_{1}\right)}}\left|\frac{f\left(\rho_{\chi_{1}}\right)}{\rho_{\chi_{1}}}\right|=o\left(T^{\alpha}\right),\,\sum_{\underset{{\scriptstyle \gamma_{\chi_{1}}\leq T}}{\rho_{\chi_{1}}\in Z_{0}^{+}\left(\chi_{1}\right)}}\left|\frac{f\left(\overline{\rho_{\chi_{1}}}\right)}{\overline{\rho_{\chi_{1}}}}\right|=o\left(T^{\alpha}\right)
\]
\[
\sum_{\underset{{\scriptstyle \gamma_{\chi_{2}}\leq T}}{\rho_{\chi_{2}}\in Z_{0}^{+}\left(\chi_{2}\right)}}\left|\frac{g\left(\rho_{\chi_{2}}\right)}{\rho_{\chi_{2}}}\right|=o\left(T^{\alpha}\right),\,\sum_{\underset{{\scriptstyle \gamma_{\chi_{2}}\leq T}}{\rho_{\chi_{2}}\in Z_{0}^{+}\left(\chi_{2}\right)}}\left|\frac{g\left(\overline{\rho_{\chi_{2}}}\right)}{\overline{\rho_{\chi_{2}}}}\right|=o\left(T^{\alpha}\right)
\]
as $T\rightarrow+\infty$ for every $\alpha>0$. Then, the double
series
\[
\sum_{\rho_{\chi_{1}}}f\left(\rho_{\chi_{1}}\right)\sum_{\rho_{\chi_{2}}}g\left(\rho_{\chi_{2}}\right)\frac{\Gamma\left(\rho_{\chi_{1}}\right)\Gamma\left(\rho_{\chi_{2}}\right)}{\Gamma\left(\rho_{\chi_{1}}+\rho_{\chi_{2}}+k+1\right)}
\]
converges absolutely for $k\in\mathbb{R},\,k>1/2$.
\end{thm}

\begin{proof}
Fix $\chi_{1}\mod q_{1},\,\chi_{2}\mod q_{2}$. We may assume that
$\gamma_{\chi_{j}}\neq0,\,j=1,2$ because, if one them is equal to
zero, we have
\[
\sum_{\rho_{\chi_{1}}}f\left(\rho_{\chi_{1}}\right)\sum_{\rho_{\chi_{2}}:\gamma_{\chi_{2}}=0}g\left(\rho_{\chi_{2}}\right)\frac{\Gamma\left(\rho_{\chi_{1}}\right)\Gamma\left(\rho_{\chi_{2}}\right)}{\Gamma\left(\rho_{\chi_{1}}+\rho_{\chi_{2}}+k+1\right)}
\]
\[
\ll_{f}N\left(2,\chi_{2}\right)\sum_{\rho_{\chi_{1}}}\left|f\left(\rho_{\chi_{1}}\right)\right|\left|\frac{\Gamma\left(\rho_{\chi_{1}}\right)}{\Gamma\left(\rho_{\chi_{1}}+k+\frac{3}{2}\right)}\right|
\]
and the series, for $k>1/2$, trivially converges (actually it converges
for smaller $k$, but it is not important for our aims) by the Stirling's
formula
\begin{equation}
\left|\Gamma\left(x+iy\right)\right|\sim e^{-\pi\left|y\right|/2}\left|y\right|^{x-1/2},\,x_{1}<x<x_{2},\,\left|y\right|\rightarrow+\infty\label{eq:stirling}
\end{equation}
(see e.g., \cite{TIT}, section 4.42), and the same argument holds
in the case $\sum_{\rho_{\chi_{1}}:\gamma_{\chi_{1}}=0}\sum_{\rho_{\chi_{2}}}$,
hence we can assume $\gamma_{\chi_{j}}\neq0,\,j=1,2$. We start with
the case $\gamma_{\chi_{1}}>0,\,\gamma_{\chi_{2}}>0$. By the Stirling's
formula (\ref{eq:stirling}), we have
\[
\sum_{\rho_{\chi_{1}}\in Z^{+}\left(\chi_{1}\right)}\left|f\left(\rho_{\chi_{1}}\right)\right|\sum_{\rho_{\chi_{2}}\in Z^{+}\left(\chi_{2}\right)}\left|g\left(\rho_{\chi_{2}}\right)\right|\left|\frac{\Gamma\left(\rho_{\chi_{1}}\right)\Gamma\left(\rho_{\chi_{2}}\right)}{\Gamma\left(\rho_{\chi_{1}}+\rho_{\chi_{2}}+k+1\right)}\right|
\]
\[
\ll\sum_{\rho_{\chi_{1}}\in Z^{+}\left(\chi_{1}\right)}\left|f\left(\rho_{\chi_{1}}\right)\right|\sum_{\rho_{\chi_{2}}\in Z^{+}\left(\chi_{2}\right)}\left|g\left(\rho_{\chi_{2}}\right)\right|\frac{1}{\left(\gamma_{\chi_{1}}+\gamma_{\chi_{2}}\right)^{3/2+k}}
\]
\[
\ll\sum_{\rho_{\chi_{1}}\in Z^{+}\left(\chi_{1}\right)}\frac{\left|f\left(\rho_{\chi_{1}}\right)\right|}{\gamma_{\chi_{1}}^{3/4+k/2}}\sum_{\rho_{\chi_{2}}\in Z^{+}\left(\chi_{2}\right)}\frac{\left|g\left(\rho_{\chi_{2}}\right)\right|}{\gamma_{\chi_{2}}^{3/4+k/2}}
\]
by arithmetic-geometric mean, and both series converges absolutely
if $k>1/2$ since, by partial summation, taking $T>0$, we have
\[
\sum_{\rho_{\chi_{1}}:0<\gamma_{\chi_{1}}\leq T}\frac{\left|f\left(\rho_{\chi_{1}}\right)\right|}{\gamma_{\chi_{1}}^{3/4+k/2}}\ll\sum_{\rho_{\chi_{1}}:0<\gamma_{\chi_{1}}\leq T}\left|\frac{f\left(\rho_{\chi_{1}}\right)}{\rho_{\chi_{1}}}\right|\frac{1}{\gamma_{\chi_{1}}^{k/2-1/4}}
\]
\begin{equation}
=\sum_{\rho_{\chi_{1}}:0<\gamma_{\chi_{1}}\leq T}\left|\frac{f\left(\rho_{\chi_{1}}\right)}{\rho_{\chi_{1}}}\right|\frac{1}{T^{k/2-1/4}}+\frac{2k-1}{4}\int_{\delta_{\chi_{1}}}^{T}\sum_{\rho_{\chi_{1}}:0<\gamma_{\chi_{1}}\leq t}\left|\frac{f\left(\rho_{\chi_{1}}\right)}{\rho_{\chi_{1}}}\right|t^{-k/2-3/4}dt\label{eq:PS double series}
\end{equation}
where $\delta_{\chi_{1}}>0$ is a number such that
\[
0<\delta_{\chi_{1}}\leq\min\left\{ \gamma_{\chi_{1}}>0:L\left(\frac{1}{2}+i\gamma_{\chi_{1}},\chi_{1}\right)=0\right\} .
\]
and the same considerations holds for the case $g\left(\rho_{\chi_{2}}\right)$.
Now, we consider the case $\gamma_{\chi_{1}}>0,\,\gamma_{\chi_{2}}<0.$
Clearly, if $\rho_{\chi_{2}}=1/2+i\gamma_{\chi_{2}}$ is a zero of
$L\left(s,\chi_{2}\right)$, then $\overline{\rho_{\chi_{2}}}$ is
a zero of $L\left(s,\overline{\chi_{2}}\right)$. Let us define
\[
\rho_{\chi_{2}}^{*}:=\overline{\rho_{\overline{\chi_{2}}}}.
\]
Then, we have to deal with
\[
\sum_{\rho_{\chi_{1}}\in Z^{+}\left(\chi_{1}\right)}\left|f\left(\rho_{\chi_{1}}\right)\right|\sum_{\rho_{\overline{\chi_{2}}}\in Z^{+}\left(\overline{\chi_{2}}\right)}\left|g\left(\rho_{\chi_{2}}^{*}\right)\right|\left|\frac{\Gamma\left(\rho_{\chi_{1}}\right)\Gamma\left(\rho_{\chi_{2}}^{*}\right)}{\Gamma\left(\rho_{\chi_{1}}+\rho_{\chi_{2}}^{*}+k+1\right)}\right|.
\]
Let $0<\alpha<1$ be fixed. We split the second series in the following
way:
\[
\sum_{\rho_{\chi_{2}}\in Z^{+}\left(\overline{\chi_{2}}\right)}=\sum_{\underset{{\scriptstyle \left|\gamma_{\chi_{1}}-\gamma_{\chi_{2}}\right|<\alpha\max\left(\gamma_{\chi_{1}},\gamma_{\chi_{2}}\right)}}{\rho_{\overline{\chi_{2}}}\in Z^{+}\left(\overline{\chi_{2}}\right)}}+\sum_{\underset{{\scriptstyle \left|\gamma_{\chi_{1}}-\gamma_{\chi_{2}}\right|\geq\alpha\max\left(\gamma_{\chi_{1}},\gamma_{\chi_{2}}\right)}}{\rho_{\overline{\chi_{2}}}\in Z^{+}\left(\overline{\chi_{2}}\right)}}
\]
\[
=:\sideset{}{_{1}}\sum+\sideset{}{_{2}}\sum,
\]
say. We start with $\sideset{}{_{1}}\sum$. From the relation
\[
\left|\Gamma\left(x+iy\right)\right|\geq\Gamma\left(x\right)\text{sech}\left(\pi y\right)^{1/2},\,x\geq1/2
\]

(see \cite{O2010}, relation $5.6.7$) we have
\[
\left|\Gamma\left(\rho_{\chi_{1}}+\rho_{\chi_{2}}^{*}+k+1\right)\right|\geq\Gamma\left(2+k\right)\min_{0\leq t\leq\alpha\max\left(\gamma_{\chi_{1}},\gamma_{\chi_{2}}\right)}\text{sech}\left(\pi t\right)^{1/2}
\]
\[
\gg_{k}e^{-\alpha\pi\left(\gamma_{\chi_{1}}+\gamma_{\chi_{2}}\right)/2}
\]
hence, using again (\ref{eq:stirling}), we get
\[
\sum_{\rho_{\chi_{1}}\in Z^{+}\left(\chi_{1}\right)}\left|f\left(\rho_{\chi_{1}}\right)\right|\sum_{\underset{{\scriptstyle \left|\gamma_{\chi_{1}}-\gamma_{\chi_{2}}\right|<\alpha\max\left(\gamma_{\chi_{1}},\gamma_{\chi_{2}}\right)}}{\rho_{\overline{\chi_{2}}}\in Z^{+}\left(\overline{\chi_{2}}\right)}}\left|g\left(\rho_{\chi_{2}}^{*}\right)\right|\left|\frac{\Gamma\left(\rho_{\chi_{1}}\right)\Gamma\left(\rho_{\chi_{2}}^{*}\right)}{\Gamma\left(\rho_{\chi_{1}}+\rho_{\chi_{2}}^{*}+k+1\right)}\right|
\]
\[
\ll_{k}\sum_{\rho_{\chi_{1}}\in Z^{+}\left(\chi_{1}\right)}\left|f\left(\rho_{\chi_{1}}\right)\right|e^{-\pi\left(1-\alpha\right)\gamma_{\chi_{1}}/2}\sum_{\rho_{\overline{\chi_{2}}}\in Z^{+}\left(\overline{\chi_{2}}\right)}\left|g\left(\rho_{\chi_{2}}^{*}\right)\right|e^{-\pi\left(1-\alpha\right)\gamma_{\chi_{2}}/2}
\]
and the convergence follows again by partial summation. Now we consider
$\sideset{}{_{2}}\sum$. We split in two further cases: if $\left|\gamma_{\chi_{1}}-\gamma_{\chi_{2}}\right|\geq\alpha\max\left(\gamma_{\chi_{1}},\gamma_{\chi_{2}}\right)$,
then $\gamma_{\chi_{1}}\neq\gamma_{\chi_{2}}$, so if $\gamma_{\chi_{1}}>\gamma_{\chi_{2}}$,
then again by (\ref{eq:stirling})
\[
\sum_{\rho_{\chi_{1}}\in Z^{+}\left(\chi_{1}\right)}\left|f\left(\rho_{\chi_{1}}\right)\right|\sum_{\underset{{\scriptstyle \left|\gamma_{\chi_{1}}-\gamma_{\chi_{2}}\right|\geq\alpha\gamma_{\chi_{1}}}}{\rho_{\overline{\chi_{2}}}\in Z^{+}\left(\overline{\chi_{2}}\right)}}\left|g\left(\rho_{\chi_{2}}^{*}\right)\right|\left|\frac{\Gamma\left(\rho_{\chi_{1}}\right)\Gamma\left(\rho_{\chi_{2}}^{*}\right)}{\Gamma\left(\rho_{\chi_{1}}+\rho_{\chi_{2}}^{*}+k+1\right)}\right|
\]
\[
\ll\sum_{\rho_{\chi_{1}}\in Z^{+}\left(\chi_{1}\right)}\left|f\left(\rho_{\chi_{1}}\right)\right|\sum_{\underset{{\scriptstyle \left|\gamma_{\chi_{1}}-\gamma_{\chi_{2}}\right|\geq\alpha\gamma_{\chi_{1}}}}{\rho_{\overline{\chi_{2}}}\in Z^{+}\left(\overline{\chi_{2}}\right)}}\left|g\left(\rho_{\chi_{2}}^{*}\right)\right|\frac{e^{-\pi\gamma_{\chi_{2}}/2}e^{-\pi\gamma_{\chi_{1}}/2}}{\left(\gamma_{\chi_{1}}-\gamma_{\chi_{2}}\right)^{3/2+k}e^{-\pi\left(\gamma_{\chi_{1}}-\gamma_{\chi_{2}}\right)/2}}
\]
\[
\ll_{\alpha}\sum_{\rho_{\chi_{1}}\in Z^{+}\left(\chi_{1}\right)}\frac{\left|f\left(\rho_{\chi_{1}}\right)\right|}{\gamma_{\chi_{1}}^{3/2+k}}\sum_{\rho_{\overline{\chi_{2}}}\in Z^{+}\left(\overline{\chi_{2}}\right)}\left|g\left(\rho_{\chi_{2}}^{*}\right)\right|e^{-\pi\gamma_{\chi_{2}}}
\]
and again the convergence follows. A similar calculation holds in
the case $\gamma_{\chi_{2}}>\gamma_{\chi_{1}}$. 
\end{proof}

We also recall the following very easy fact, because we will use it
very often.
\begin{prop}
\label{prop:trivial}For every $a>0,\,\alpha\geq0$ and $n\geq0$
we have
\[
\int_{0}^{a}x^{\alpha}\left|\log\left(x\right)\right|^{n}dx\ll_{n}a^{\alpha+1}\left(1+\left|\log\left(a\right)\right|^{n}\right).
\]
\end{prop}

\begin{thebibliography}{99}
\bibitem[B2011]{B2011} H. Brezis. Functional analysis, Sobolev spaces and partial differential equations, Vol. 2. No. 3, New York: Springer, 2011.
\bibitem[BP2018]{BP2018} M. B\u{a}nescu, D. Popa, A multiple Abel summation formula and asymptotic evaluations for multiple sums, International Journal of Number Theory, 14(4) (2018), 1197--1210.
\bibitem[C1995]{C1995} J.-Y. Chemin, Fluides parfaits incompressibles, Ast\'erisque, 230. Soci\'et\'e Math\'ematique de France, Paris, 1995.
\bibitem[CGZ2025]{CGZ2025} M. Cantarini, A. Gambini, A. Zaccagnini, Laplace convolutions of weighted averages of arithmetical functions, Forum Mathematicum 37(2) (2025), 515--533.
\bibitem[CGZ2026]{CGZ2026} M. Cantarini, A. Gambini, A. Zaccagnini, On the discrete convolution of the Liouville and M\"obius function, submitted. https://arxiv.org/abs/2603.10241
\bibitem[DAVEN]{DAVEN} H. Davenport, Multiplicative Number Theory, volume 74 of Graduate Texts in Mathematics. Springer-Verlag, New York, third edition, 2000. Revised and with a preface by Hugh L. Montgomery.
\bibitem[DVL1993]{DVL1993} R. A. DeVore, G. G. Lorentz, Constructive approximation, Vol. 303, Springer Science \& Business Media, 1993.
\bibitem[K1983]{K1983} S. G. Krantz, Lipschitz spaces, smoothness of functions, and approximation theory, Exposition. Math. 3 (1983), 193--260.
\bibitem[KT2006]{KT2006} I. Kiuchi, Y. Tanigawa,Bounds for double zeta-functions, Annali della Scuola Normale Superiore di Pisa-Classe di Scienze, 5(4) (2006), 445--464.
\bibitem[MV2006]{MV2006} H. L. Montgomery, R. C. Vaughan, Multiplicative number theory I: Classical theory. Vol. 97. Cambridge university press, 2006.
\bibitem[O2010]{O2010} F. W. J. Olver, ed. NIST handbook of mathematical functions hardback and CD-ROM. Cambridge university press, 2010.
\bibitem[R2022]{R2022} A. Rainer, H\"older--Zygmund classes on smooth curves. Zeitschrift f\"ur Analysis und ihre Anwendungen, 41(1-2) (2022, 189--209.
\bibitem[TIT]{TIT} E. C. Titchmarsh, The Theory of Functions, 2nd ed., Oxford Univ. Press, 1988.
\end{thebibliography}
\end{document}
